\section{Contexto y Fronteras del Sistema}

En línea con el enfoque de la Teoría General de Sistemas, es fundamental definir el contexto y las fronteras que delimitan el alcance del software a construir.

\begin{itemize}
    \item \textbf{Contexto:} El sistema opera dentro del contexto más amplio de la \textbf{gestión de activos y operaciones de una planta industrial}. Su propósito es interactuar con los procesos de mantenimiento para optimizar la fiabilidad y productividad de los equipos.
    \item \textbf{Fronteras:} Las fronteras del sistema se definen por sus puntos de interacción:
    \begin{itemize}
        \item \textbf{Humanas:} Los \textbf{Actores Responsables} (Técnicos, Supervisores, Gerentes, etc.) son las interfaces humanas que operan dentro del sistema o interactúan con él.
        \item \textbf{Técnicas:} Los \textbf{Procesos de Interoperabilidad} definen las fronteras con otros sistemas de software. El sistema se conectará con plataformas ERP/EAM existentes y con herramientas de productividad (Google Workspace/Office 365), pero no es responsable de las funcionalidades internas de dichos sistemas. Su límite está en el intercambio de datos con ellos.
    \end{itemize}
\end{itemize}

